\subsection{Rancangan Alur Akses dan Delegasi}
\label{subsec:rancangan-alur-akses-dan-delegasi}

Fitur \textit{offline attenuation} dirancang untuk menjawab kebutuhan delegasi akses antar tenaga medis tanpa mengharuskan pasien menerbitkan ulang token setiap kali terjadi perpindahan tugas. Pada rancangan ini, istilah \textit{offline} tidak berarti bahwa seluruh proses akses data dilakukan tanpa koneksi ke \textit{Server} PRE. Istilah tersebut merujuk pada kemampuan pemegang token untuk menurunkan atau mengatenuasi token secara lokal pada aplikasi \textit{client}, tanpa meminta \textit{Server} PRE menerbitkan token baru.

Delegasi akses pada rancangan ini terdiri atas dua tahap utama. Tahap pertama adalah delegasi awal dari pasien kepada dokter penanggung jawab pasien. Pada tahap ini, pasien memberikan persetujuan akses, kemudian \textit{Server} PRE menerbitkan token Macaroon awal untuk dokter. Tahap kedua adalah delegasi lanjutan dari dokter kepada aktor layanan kesehatan lain, seperti perawat, petugas laboratorium, atau apoteker. Pada tahap kedua ini, token tidak diterbitkan ulang oleh \textit{Server} PRE, melainkan diturunkan secara lokal oleh aplikasi \textit{client} dokter dengan menambahkan \textit{caveats} baru yang lebih ketat.

Macaroons mendukung atenuasi lokal karena token turunan dapat dibuat dengan menambahkan \textit{caveats} baru pada token induk. Penambahan \textit{caveats} menghasilkan tanda tangan Macaroon baru melalui mekanisme \textit{chained-HMAC}. Karena \textit{caveats} pada Macaroon bersifat \textit{append-only}, token turunan tidak dapat menghapus batasan pada token induk. Dengan demikian, token turunan hanya dapat memiliki cakupan akses yang sama atau lebih sempit daripada token induknya.

\subsubsection{Delegasi Awal dari Pasien kepada Dokter}

Delegasi awal merupakan proses pemberian otorisasi pertama dari pasien kepada dokter penanggung jawab pasien. Pada tahap ini, pasien tetap menjadi pihak yang mengendalikan pemberian akses awal terhadap data RME miliknya. Namun, karena \texttt{RootKey} Macaroon hanya dimiliki oleh \textit{Server} PRE, pasien tidak membuat token Macaroon secara langsung. Pasien memberikan persetujuan melalui \textit{Patient Client}, kemudian \textit{Server} PRE menerbitkan token awal berdasarkan persetujuan tersebut.

Alur delegasi awal dari pasien kepada dokter adalah sebagai berikut.

\begin{enumerate}[leftmargin=*]
	\item Pasien memilih konteks jenis pelayanan kesehatan dan melakukan scan kode QR yang berisi data	\texttt{iota\_address} dan \texttt{pre\_public\_key} dari dokter penanggung jawab pasien.
	\item Pasien menentukan konteks layanan kesehatan, misalnya rawat jalan.
	\item \textit{Patient Client} mengirimkan permintaan penerbitan token kepada \textit{Server} PRE.
	\item \textit{Server} PRE memvalidasi identitas pasien dan cakupan akses yang diminta.
	\item \textit{Server} PRE menerbitkan Macaroon awal dengan \textit{caveats} yang mengikat token pada pasien, RME terkait, dokter penerima akses, cakupan baca dan tulis, masa berlaku, serta batas kedalaman delegasi.
	\item Token awal diberikan kepada dokter melalui \textit{Hospital Client} atau disimpan pada penyimpanan token aplikasi yang hanya dapat diakses oleh dokter terkait.
\end{enumerate}

Contoh \textit{caveats} pada token awal yang diberikan kepada dokter dalam konteks pelayanan rawat jalanditunjukkan pada Listing \ref{lst:caveat-token-awal-dokter}.

\begin{lstlisting}[caption={Contoh \textit{Caveats} Token Awal dari Pasien kepada Dokter}, label={lst:caveat-token-awal-dokter}]
patient_address = 0xPASIEN
related_rme_id = RME-001
root_subject = 0xDOKTER

read_dataset_in = [RAWAT_JALAN, LABORATORIUM, APOTEK]
write_dataset_in = [RAWAT_JALAN, LABORATORIUM, APOTEK]

read_function_in = [
  ANAMNESIS,
  PEMERIKSAAN_FISIK,
  PEMERIKSAAN_PSIKOLOGIS,
  DIAGNOSIS,
  TERAPI,
  PERMINTAAN_PEMERIKSAAN,
  HASIL_PEMERIKSAAN,
  DATA_RESEP_DAN_OBAT
]

write_function_in = [
  ANAMNESIS,
  PEMERIKSAAN_FISIK,
  DIAGNOSIS,
  TERAPI,
  PERMINTAAN_PEMERIKSAAN
]

expires_before = 2026-05-16T18:00:00
max_delegation_depth = 1
\end{lstlisting}

Pada contoh tersebut, \texttt{root\_subject} menyatakan bahwa dokter merupakan penerima token awal. \texttt{patient\_address} mengikat token pada data pasien tertentu, sedangkan \texttt{related\_rme\_id} membatasi token pada episode RME tertentu. \texttt{read\_dataset\_in}, \texttt{write\_dataset\_in}, \texttt{read\_function\_in}, dan \texttt{write\_function\_in} menyatakan cakupan akses yang diberikan. \texttt{expires\_before} membatasi masa berlaku token. Sementara itu, \texttt{max\_delegation\_depth} membatasi jumlah delegasi lanjutan yang dapat dilakukan dari token tersebut.

\subsubsection{Delegasi Lokal dari Dokter kepada Aktor Lain}

Setelah dokter menerima token awal, dokter dapat menurunkan token tersebut kepada aktor layanan kesehatan lain apabila dibutuhkan dalam proses pelayanan. Delegasi ini dilakukan secara lokal pada aplikasi \textit{client} dokter dengan menambahkan \textit{caveats} baru pada token induk. \textit{Server} PRE tidak menerbitkan token baru pada tahap ini.

Aktor yang dapat menerima token turunan dari dokter mencakup perawat, petugas laboratorium, dan apoteker. Setiap aktor hanya menerima cakupan akses sesuai kebutuhan tugasnya. Contoh pemetaan kebutuhan delegasi ditunjukkan pada Tabel \ref{tab:pemetaan-delegasi-dokter}.

\begin{table}[!ht]
	\centering
	\caption{Pemetaan Delegasi Token dari Dokter kepada Aktor Lain}
	\label{tab:pemetaan-delegasi-dokter}
	\small
	\begin{tabular}{|p{0.18\textwidth}|p{0.23\textwidth}|p{0.43\textwidth}|}
		\hline
		\textbf{Penerima Delegasi} & \textbf{Tujuan Delegasi}                                                                                                                                             & \textbf{Cakupan Akses} \\ \hline

		Perawat
		                           & Membantu pengisian dan pembacaan data klinis tertentu.
		                           & Membaca dan/atau menulis segmen seperti anamnesis, pemeriksaan fisik, pemeriksaan psikologis, serta instruksi medik dan keperawatan sesuai kebutuhan layanan.
		\\ \hline

		Petugas laboratorium
		                           & Memproses permintaan pemeriksaan dan mengunggah hasil laboratorium.
		                           & Membaca segmen permintaan pemeriksaan dan menulis segmen hasil pemeriksaan pada \textit{dataset} laboratorium.
		\\ \hline

		Apoteker
		                           & Memproses resep dan penyerahan obat.
		                           & Membaca segmen resep, alergi, dan informasi obat yang relevan, serta menulis status pengkajian resep, status resep, waktu penyiapan obat, dan waktu penyerahan obat.
		\\ \hline
	\end{tabular}
\end{table}

Sebagai contoh, apabila dokter membutuhkan pemeriksaan laboratorium, aplikasi \textit{client} dokter membuat token turunan untuk petugas laboratorium dengan menambahkan \textit{caveats} baru seperti pada Listing \ref{lst:caveat-token-turunan-lab}.

\begin{lstlisting}[caption={Contoh \textit{Caveats} Tambahan Token Turunan untuk Petugas Laboratorium}, label={lst:caveat-token-turunan-lab}]
delegated_by = 0xDOKTER
delegated_to = 0xLAB

read_dataset_in = [LABORATORIUM]
write_dataset_in = [LABORATORIUM]

read_function_in = [PERMINTAAN_PEMERIKSAAN]
write_function_in = [HASIL_PEMERIKSAAN]

expires_before = 2026-05-16T14:00:00
max_delegation_depth = 1
\end{lstlisting}

Pada contoh tersebut, token turunan hanya dapat digunakan oleh petugas laboratorium untuk membaca permintaan pemeriksaan dan menulis hasil pemeriksaan. Token tidak dapat digunakan untuk membaca diagnosis, anamnesis, terapi, resep, atau segmen klinis lain yang tidak termasuk dalam \textit{caveats} token turunan.

Nilai \texttt{max\_delegation\_depth} pada rancangan ini diperlakukan sebagai batas maksimum panjang rantai delegasi, bukan sebagai nilai yang ditimpa secara bebas. \textit{Server} PRE menghitung jumlah \textit{caveat} \texttt{delegated\_to} pada token untuk menentukan kedalaman delegasi aktual. Apabila token awal memiliki \texttt{max\_delegation\_depth = 1}, maka token tersebut hanya boleh diturunkan satu kali dari dokter kepada aktor lain. Jika aktor penerima mencoba menurunkan token lagi kepada pihak berikutnya, jumlah delegasi akan melebihi batas dan token ditolak oleh \textit{Server} PRE.

\subsubsection{Proses Pemberian Token Turunan}

Proses pemberian token turunan dari dokter kepada aktor lain tidak dilakukan dengan meminta \textit{Server} PRE membuat token baru. Dokter menurunkan token secara lokal, kemudian token turunan diberikan kepada aktor penerima melalui kanal aplikasi yang terautentikasi, misalnya melalui \textit{Hospital Client} atau penyimpanan token internal yang dikelola oleh sistem.

Agar proses delegasi dapat diaudit dan tidak bergantung hanya pada isi token, aplikasi dokter juga membuat bukti delegasi yang ditandatangani menggunakan kunci privat \textit{wallet} dokter. Bukti ini digunakan untuk menunjukkan bahwa token turunan benar-benar dibuat oleh pihak yang dinyatakan pada \texttt{delegated\_by}. Bukti delegasi dapat disertakan sebagai metadata pendamping token atau dicatat sebagai \textit{event} pada sistem.

Struktur konseptual pemberian token turunan ditunjukkan pada Listing \ref{lst:delegation-envelope}.

\begin{lstlisting}[caption={Struktur Konseptual Pemberian Token Turunan}, label={lst:delegation-envelope}]
DelegationEnvelope {
  macaroon_token: "<attenuated_macaroon>",
  delegated_by: "0xDOKTER",
  delegated_to: "0xLAB",
  related_rme_id: "RME-001",
  issued_at: "2026-05-16T10:30:00",
  expires_before: "2026-05-16T14:00:00",
  token_hash: "hash(attenuated_macaroon)",
  delegation_signature: "signature_by_0xDOKTER"
}
\end{lstlisting}

\texttt{delegation\_signature} merupakan tanda tangan digital dari dokter terhadap konteks delegasi, misalnya \texttt{delegated\_to}, \texttt{related\_rme\_id}, \texttt{expires\_before}, dan \texttt{token\_hash}. Dengan demikian, pihak lain yang hanya memperoleh token dokter secara tidak sah tidak dapat membuat delegasi yang valid atas nama dokter tanpa memiliki kunci privat \textit{wallet} dokter.

Secara umum, proses pemberian token turunan berlangsung sebagai berikut.

\begin{enumerate}[leftmargin=*]
	\item Dokter memilih aktor penerima delegasi melalui \textit{Hospital Client}, misalnya petugas laboratorium.
	\item Aplikasi dokter menentukan cakupan akses yang akan diberikan, seperti \texttt{LABORATORIUM} dan \texttt{HASIL\_PEMERIKSAAN}.
	\item Aplikasi dokter menambahkan \textit{caveats} baru ke Macaroon induk sehingga terbentuk Macaroon turunan.
	\item Aplikasi dokter menghitung \textit{hash} dari Macaroon turunan.
	\item Dokter menandatangani konteks delegasi menggunakan kunci privat \textit{wallet}.
	\item Macaroon turunan beserta bukti delegasi dikirimkan kepada aktor penerima melalui kanal aplikasi yang terautentikasi.
	\item Aktor penerima menyimpan token turunan dan menggunakannya ketika membutuhkan akses ke segmen RME terkait.
\end{enumerate}

\subsubsection{Contoh Delegasi kepada Perawat, Laboratorium, dan Apotek}

Contoh \textit{caveats} tambahan untuk token turunan kepada perawat ditunjukkan pada Listing \ref{lst:caveat-token-turunan-perawat}.

\begin{lstlisting}[caption={Contoh \textit{Caveats} Tambahan Token Turunan untuk Perawat}, label={lst:caveat-token-turunan-perawat}]
delegated_by = 0xDOKTER
delegated_to = 0xPERAWAT

read_dataset_in = [RAWAT_JALAN]
write_dataset_in = [RAWAT_JALAN]

read_function_in = [
  ANAMNESIS,
  PEMERIKSAAN_FISIK,
  PEMERIKSAAN_PSIKOLOGIS,
  INSTRUKSI_MEDIK_DAN_KEPERAWATAN
]

write_function_in = [
  ANAMNESIS,
  PEMERIKSAAN_FISIK,
  PEMERIKSAAN_PSIKOLOGIS
]

expires_before = 2026-05-16T16:00:00
max_delegation_depth = 1
\end{lstlisting}

Contoh \textit{caveats} tambahan untuk token turunan kepada apoteker ditunjukkan pada Listing \ref{lst:caveat-token-turunan-apoteker}.

\begin{lstlisting}[caption={Contoh \textit{Caveats} Tambahan Token Turunan untuk Apoteker}, label={lst:caveat-token-turunan-apoteker}]
delegated_by = 0xDOKTER
delegated_to = 0xAPOTEKER

read_dataset_in = [APOTEK]
write_dataset_in = [APOTEK]

read_function_in = [
  DATA_RESEP_DAN_OBAT,
  RIWAYAT_ALERGI,
  ASAL_RESEP,
  DOKTER_PENULIS_RESEP
]

write_function_in = [
  STATUS_DAN_PENGKAJIAN_RESEP,
  STATUS_RESEP,
  WAKTU_PENYIAPAN_OBAT,
  WAKTU_PENYERAHAN_OBAT,
  PETUGAS_DISPENSING
]

expires_before = 2026-05-16T17:00:00
max_delegation_depth = 1
\end{lstlisting}

Pada kedua contoh tersebut, token turunan tetap membawa seluruh \textit{caveats} dari token dokter. Oleh karena itu, apabila token dokter hanya berlaku untuk \texttt{related\_rme\_id = RME-001}, maka token perawat, token laboratorium, dan token apoteker juga hanya berlaku untuk RME tersebut. Apabila token dokter berlaku sampai pukul 18.00, token turunan tidak dapat memperpanjang masa berlaku melebihi waktu tersebut karena \textit{Server} PRE akan memvalidasi seluruh \textit{caveats} \texttt{expires\_before} dan menggunakan batas waktu yang paling ketat.

\subsubsection{Verifikasi Token Turunan oleh \textit{Server} PRE}

Ketika aktor penerima delegasi menggunakan token turunan, \textit{Server} PRE bertindak sebagai titik penegakan kebijakan akses. \textit{Server} PRE tidak menerbitkan ulang token, tetapi memverifikasi token turunan dan seluruh bukti yang menyertainya.

Langkah verifikasi token turunan adalah sebagai berikut.

\begin{enumerate}[leftmargin=*]
	\item \textit{Server} PRE memverifikasi validitas Macaroon menggunakan \texttt{RootKey}.
	\item \textit{Server} PRE memverifikasi bahwa \texttt{patient\_address} dan \texttt{related\_rme\_id} pada token sesuai dengan data RME yang diminta.
	\item \textit{Server} PRE menghitung aktor aktif. Jika tidak terdapat \texttt{delegated\_to}, maka aktor aktif adalah \texttt{root\_subject}. Jika terdapat satu atau lebih \texttt{delegated\_to}, maka aktor aktif adalah nilai \texttt{delegated\_to} terakhir.
	\item \textit{Server} PRE memverifikasi tanda tangan \textit{wallet} dari aktor aktif pada permintaan akses.
	\item \textit{Server} PRE memverifikasi bukti delegasi dari pihak yang dinyatakan pada \texttt{delegated\_by}.
	\item \textit{Server} PRE menghitung kedalaman delegasi berdasarkan jumlah \texttt{delegated\_to} pada token dan memastikan nilainya tidak melebihi \texttt{max\_delegation\_depth}.
	\item \textit{Server} PRE memvalidasi seluruh batas waktu pada \texttt{expires\_before} dan menggunakan batas waktu yang paling awal sebagai masa berlaku efektif token.
	\item \textit{Server} PRE menghitung cakupan akses efektif berdasarkan irisan seluruh \texttt{read\_dataset\_in}, \texttt{write\_dataset\_in}, \texttt{read\_function\_in}, dan \texttt{write\_function\_in}.
	\item \textit{Server} PRE mengambil metadata segmen dari IOTA dan mencocokkan \texttt{dataset\_category} serta \texttt{function\_category} segmen dengan cakupan akses efektif token.
	\item Jika seluruh verifikasi berhasil, \textit{Server} PRE melanjutkan proses pengambilan data terenkripsi dari IPFS dan melakukan \textit{re-encryption}. Jika salah satu pemeriksaan gagal, permintaan akses ditolak.
\end{enumerate}

Dengan mekanisme tersebut, token turunan tidak dapat memperluas hak akses token induk. Misalnya, apabila token dokter tidak memiliki akses terhadap \texttt{APOTEK}, maka dokter tidak dapat membuat token turunan yang valid untuk apoteker. Meskipun dokter menambahkan \textit{caveat} \texttt{read\_dataset\_in = [APOTEK]}, cakupan efektif token akan menjadi kosong karena batasan tersebut tidak termasuk dalam cakupan token induk. Demikian pula, apabila token dokter hanya berlaku hingga pukul 18.00, token turunan yang mencoba menggunakan \texttt{expires\_before} setelah pukul 18.00 tetap tidak akan memperpanjang masa berlaku efektif token.

Rancangan ini mempertahankan kendali awal pasien karena token awal tetap diterbitkan berdasarkan persetujuan pasien. Pada saat yang sama, rancangan ini mengurangi hambatan operasional pada pelayanan klinis karena dokter dapat menurunkan token secara lokal kepada aktor lain yang terlibat dalam layanan. Dengan demikian, mekanisme delegasi ini mendukung kolaborasi layanan kesehatan yang dinamis, menjaga prinsip \textit{least privilege}, dan tetap mempertahankan peran \textit{Server} PRE sebagai titik verifikasi dan penegakan akses.